\chapter*{Appendice C --- Mappe e Handout}
\addcontentsline{toc}{chapter}{Appendice C --- Mappe e Handout}
\markboth{\MakeUppercase{Appendice C --- Mappe e Handout}}{}
\label{chap:appendiceC}

\begin{DndComment}[]{Nota editoriale}
  Quattro handout per il Capitolo 1, pensati per essere mostrati ai giocatori nel momento esatto della scoperta (lettura ad alta voce, o reveal su Foundry/The Forge). Il testo del vero e proprio ``Canto del Seme Velato'' (\texttt{the-prophetic-song.md}) resta IPOTESI, non ancora confermato dall'autore -- l'handout sotto usa invece le parole di Fendrel già CONFERMATE nel testo di Capitolo 1, non la versione estesa del canto.
\end{DndComment}

\section{Le parole di Fendrel}

\begin{center}
  \includegraphics[width=6cm]{pergamena-profezia-bg}
\end{center}

\begin{quotation}
\noindent \textit{``Ascoltatemi, perché non avrò il tempo di ripetermi. Esiste un canto antico, custodito da questo circolo da prima che io nascessi:''}

\textit{``Nel dì in cui l'eco del cosmo si tingerà dei colori dell'antitesi, al mondo verrà colui che in sé racchiude le luci e le ombre, e i doni che farà saranno l'impronta del suo essere e del suo auspicio -- né nefasto né propizio, ma della stessa natura delle cose.''}

\textit{``Non vi chiedo di capire cosa significhi -- non lo capisco nemmeno io. Vi chiedo questo: quando il seme torna, qualcuno deve salvarlo e crescerlo, lontano da chi lo vuole usare. La profezia dice che sarete voi.''}
\end{quotation}

\textit{Uso consigliato.} Mostrare nell'istante in cui Fendrel parla (Capitolo 1, Zona 4) -- lasciare che il testo parli da solo, senza commento.

\section{La mappa incisa della Sala della Guarnigione}

\begin{center}
  \includegraphics[width=7cm]{mappa-sala-guarnigione}
\end{center}

Trovata in Capitolo 1, Zona 8 (Investigazione CD 11): mostra la disposizione delle gallerie sottostanti e il bivio di Parte 2 -- la via diretta, più vicina alla corruzione residua, e la via più lunga ma sicura.

\textit{Uso consigliato.} Mostrare nel momento in cui i giocatori devono scegliere quale via seguire -- è un supporto decisionale reale, non solo colore.

\section{Il frammento di mappa orchesco}

\begin{center}
  \includegraphics[width=7cm]{mappa-orchesca}
\end{center}

Trovato in Capitolo 1, Zona 4 di Aelothir (Deposito Nascosto, Investigazione o Percezione CD 13): segna il percorso verso il campo di Grulnak e un'area evitata, marcata con un simbolo di avvertimento senza spiegazione.

\section{La tavoletta d'avorio}

\begin{center}
  \includegraphics[width=7cm]{tavoletta-avorio}
\end{center}

Trovata in Capitolo 1, Zona 9b (Archivio Sommerso, Investigazione CD 14): un rapporto di guarnigione che nomina ``il varco inferiore'' e il Capitano Ethandrel Nightsong.
